\section{Exhausting the normalization boundary}

This section gives the codimension-one calculation on which the rest of the
paper depends.  A geometric branch over $Z$ means a prime after passing from
the generic DVR of $Z$ to its strict henselization; it has residue degree one.
The tables below are instead written over the displayed generic DVR whenever
a moving residual polynomial is more naturally kept irreducible.  After
strict henselization such a row splits into its geometric branches.  In
general the arithmetic residue degree is the degree of the corresponding
irreducible residual factor.

Let $Z\subset Y$ be a prime divisor, let
${\cal O}_{Y,\eta_Z}$ have uniformizer $\pi_Z$, and let $S_Z$ be its integral
closure in $L=k(U)$.  Normality makes $S_Z$ a finite semilocal Dedekind
domain.  If $[L:K]=N$, then
\[
 N=\dim_{k(Z)}S_Z/\pi_ZS_Z
  =\sum_{E\mid Z}e(E/Z)f(E/Z).                            \tag{4.1}
\]
Every summand is positive and the sum includes primes whose generic points
belong to the affine-source open.  Thus a displayed list which reaches $N$
is the complete list over $Z$.

We use reduced discriminant equations.  The critical-value maps below have
integral source and height-one prime kernel in the target UFD.  We denote a
generator of that kernel by $\delta$ and put $Z_\Delta=V(\delta)$.  This is
important for cancellation maps: their raw polynomial discriminant has
generic multiplicity $r$, whereas $\delta$ occurs once.

\begin{lemma}[exact weighted discriminant saturation]
\label{lem:weighted-saturation}
For
\[
 H_N(W)=W^2(1-W)(1+W^{N-3}),\qquad
 E_N(W)=H_N(W)-BCW+2AC^2,
\]
there are nonzero constants $\lambda_N,\mu_N\in\C^\times$ and an
irreducible polynomial $\delta_N\in\C[A,B,C]$ such that
\[
 \operatorname{Disc}_W(E_N)=\lambda_N C^2\delta_N,
 \qquad
 \delta_N|_{C=0}=\mu_N(B^2-8A).                         \tag{4.1a}
\]
In particular, $C\nmid\delta_N$, the residual discriminant
$V(\delta_N)$ is prime and reduced, and its generic normalization prime is a
boundary prime with $(e,f)=(2,1)$.
\end{lemma}

\begin{proof}
Work first over the generic DVR $\C(A,B)[[C]]$.  The zero of $H_N$ at the
origin has exact multiplicity two and leading coefficient one.  In the
scaled chart $W=Cu$,
\[
 C^{-2}E_N(Cu)=u^2-Bu+2A+O(C).                          \tag{4.1b}
\]
Thus the two roots specializing to zero differ by $C$ times a unit over the
generic quadratic extension, and the square of their difference contributes
exactly
$C^2(B^2-8A)$ to the root-product formula for the discriminant.  Every other
root specializes to one of the distinct nonzero simple roots of $H_N$; all
its differences from the two small roots and from the other simple roots are
units.  Consequently
\[
 v_C(\operatorname{Disc}_W E_N)=2,
 \quad
 \left.C^{-2}\operatorname{Disc}_W(E_N)\right|_{C=0}
   =\nu_N(B^2-8A)
\]
for some $\nu_N\in\C^\times$.  This proves that the forced power is exactly
$C^2$ and that the quotient is not divisible by $C$.

It remains to identify that quotient intrinsically.  On $C\ne0$ the
critical incidence is
\[
 B=\frac{H_N'(W)}{C},\qquad
 A=\frac{WH_N'(W)-H_N(W)}{2C^2}.
\]
It is integral and generically birational to its image: for
$s=H_N'(W)$ and $t=WH_N'(W)-H_N(W)$ one has $dt/ds=W$ wherever
$H_N''(W)\ne0$.  Its height-one kernel is therefore generated by an
irreducible polynomial $\delta_N$.  The generic critical root is exactly
double, so the discriminant vanishes to order one along this image.  Hence
the $C$-primitive quotient of the discriminant is a nonzero scalar multiple
of $\delta_N$, proving (4.1a), primeness, and reducedness.

At the generic point of $V(\delta_N)$, $C$ is a unit.  If
$u=E_{N,W}$ is the uniformizer of the normalized critical incidence, then
\[
 \delta_N=\varepsilon u^2,
 \qquad
 \varepsilon\in
 \widehat{\mathcal O}_{\overline X,E_\Delta}^{\times}.
\]
The critical-value map is birational, so $f=1$, and the displayed equation
gives $e=2$.  This prime cannot meet the affine-source open: the Keller map
is etale there, whereas this DVR extension is ramified.  It is therefore
exactly the claimed ramified boundary vertex.
\end{proof}

\begin{lemma}[boundary exhaustion]\label{lem:boundary}
The complete boundary-image list for the split weighted map is
$\Delta_H=V(\delta_H)$ and $C=0$.  Over $\Delta_H$ there is one boundary
prime with $(e,f)=(2,1)$; over $C=0$ there are $N-3$ geometric boundary
primes with $(e,f)=(1,1)$.

For the cancellation map the complete list is
$\Delta_{m,r}=V(\delta_{m,r})$ and $P=0$.  Over the first there is one
boundary prime with $(e,f)=(r+1,1)$; over the second there are $mr-1$
geometric boundary primes with $(e,f)=(1,1)$.  For $N\ge4$ both target
vertices occur.  In each family the discriminant is the unique vertex
receiving ramification.
\end{lemma}

\begin{proof}
We give the local calculation family by family.

\smallskip
\noindent\emph{Weighted discriminant.}
For the degreewise seed $H_N$, Lemma~\ref{lem:weighted-saturation} already
proves the exact $C^2$ saturation and identifies the reduced prime equation.
The following local description records the corresponding generic DVR and
also applies to any split seed satisfying the same simple-root hypotheses.
Write $E(W)=H(W)-BCW+cAC^2$.  On $C\ne0$, the critical incidence is
\[
 E=E_W=0,\quad
 B=\frac{H'(W)}{C},\quad
 A=\frac{WH'(W)-H(W)}{cC^2}.                             \tag{4.2}
\]
It is parameterized by $(W,C)\in\A^1\times\mathbb G_m$ and is integral.
It is generically birational: if $s=H'(W)$ and
$t=WH'(W)-H(W)$, then, where $H''(W)\ne0$, one has $dt/ds=W$.
Thus $W\in\C(s,t)$, while $C$ is already the third target coordinate.  Hence
its height-one kernel
$(\delta_H)\subset\C[A,B,C]$ is prime.  Consequently
$\Delta_H=V(\delta_H)$ is irreducible and reduced.

At its generic point $H''(W)\ne0$ and exactly one inverse root is double.
In the normalized local ring put $u=E_W=H'(W)-BC$.  The critical divisor is
$u=0$, so $u$ is an upstairs uniformizer.  Taylor expansion in the root
coordinate gives
\[
 \delta_H=\varepsilon u^2,\quad
 \varepsilon\in
 \widehat{\mathcal O}_{\overline X,E_\Delta}^{\times}.   \tag{4.3}
\]
Thus the generic DVR extension has $(e,f)=(2,1)$.  The ordinary polynomial
discriminant is a unit times $\delta_H$ at this DVR, not a higher power,
because the generic double-root discriminant is transverse.  Globally it
may also contain a power of $C$; equivalently, the primitive part of
$(\operatorname{Disc}_W E):C^\infty$ contains $\delta_H$ once.  The divisor
$C=0$ is audited separately below.  The critical prime is outside the affine
source: $E_W=-c\gamma$ makes reconstruction force $\gamma=0$, and a ramified
prime cannot lie in the etale Keller open.  After strict
henselization the other $N-2$ roots are simple affine branches; their local
uniformizer is $\delta_H$ and each has $(e,f)=(1,1)$.  Before strict
henselization, each irreducible factor of the residual simple-root quotient
gives one affine prime with $e=1$ and $f$ equal to that factor's degree; the
degrees sum to $N-2$.  Hence
\[
 2+(N-2)=N.                                               \tag{4.4}
\]

\smallskip
\noindent\emph{Weighted divisor $C=0$.}
The ideal $(C)$ is prime and reduced.  Its generic DVR has uniformizer $C$
and residue field $\C(A,B)$.  Since $E(W)|_{C=0}=H(W)$, its complete local
branch table is
\[
\begin{array}{c|c|c|c}
 \text{root chart}&\text{generic local equation}&(e,f)&\text{position}\\ \hline
 W=Cu&h_2u^2-Bu+cA+O(C)=0&(1,2)&\text{affine}\\
 W=1+O(C)&E_W(1)=-c&(1,1)&\text{affine}\\
 W=\rho_j+O(C)&E_W(\rho_j)=H'(\rho_j)\ne0&(1,1)&\text{boundary}.
\end{array}                                                \tag{4.5}
\]
In every row $C$ remains an upstairs uniformizer.  The first residual
quadratic is irreducible over $\C(A,B)$ and generically separable, giving
$f=2$.  Direct substitution in reconstruction gives nonnegative valuation
for every source coordinate in the first two rows.  In the last row,
\[
 y=\frac{W-\gamma}{C},\quad
 (W-\gamma)|_{C=0}=\rho_j+\frac{H'(\rho_j)}{c}\ne0,     \tag{4.6}
\]
so $y$ has valuation $-1$.  For $H_N$ the $\rho_j$ are the $N-3$ distinct
roots of $1+W^{N-3}$.  The displayed numerator is nonzero for these roots
when $N\ge5$.  In the sole exceptional case $N=4$, $\rho=-1$ makes its
constant term vanish, but the next normalized jet is
\[
 W=-1-BC+\frac{3B^2-A}{2}C^2+O(C^3),
 \quad C^2z\longrightarrow2,                            \tag{4.7}
\]
so that branch is still boundary.  The degree sum is
\[
 2+1+(N-3)=N.                                             \tag{4.8}
\]

\smallskip
\noindent\emph{Cancellation discriminant.}
Put $D(T)=1-T(Q-PT)^m$.  On $P\ne0$ the inverse polynomial is monic up to a
unit and $\Psi_T=CD(T)^r$.  With $Y_0=Q-PT$, the critical incidence has
\[
 T=Y_0^{-m},\quad P=(Q-Y_0)Y_0^m,                        \tag{4.9}
\]
and $R$ equal to the critical value.  It is integral, and generic critical
values are distinct.  To see the latter without a genericity assertion,
specialize $Q=0$: the $m+1$ critical roots differ by $(m+1)$-st roots of
unity and their critical values are those roots times the same nonzero beta
integral $C\int_0^1(1-s^{m+1})^r\,ds$.  Thus a target point on a dense open
of the discriminant has a unique critical preimage, so the critical-value
map is birational to its image.  Its height-one
kernel $(\delta_{m,r})\subset\C[P,Q,R]$ is prime and defines the reduced
irreducible divisor $\Delta_{m,r}$.

At the generic critical point $u=D(T)$ is an upstairs uniformizer.  Since
$D$ has a simple zero in the root direction, integration gives
\[
 \delta_{m,r}=\varepsilon u^{r+1},\quad
 \varepsilon\in
 \widehat{\mathcal O}_{\overline X,E_\Delta}^{\times}.   \tag{4.10}
\]
Hence the critical boundary prime has $(e,f)=(r+1,1)$.  The raw polynomial
discriminant is generically a unit times $\delta_{m,r}^{r}$; further powers
of $P$ arise only from clearing the leading coefficient of the $T$ chart.
The reduced equation used here is the primitive generator of
\[
 \sqrt{(\operatorname{Res}_T(\Psi,D)):P^\infty}
   =(\delta_{m,r}),                                       \tag{4.11}
\]
not the scheme discriminant.  The ramified prime is outside the affine
source because $A=D^{-1}$ has negative valuation.  All other geometric roots
are simple affine branches with $(e,f)=(1,1)$ after strict henselization.
Before that base change, the irreducible factors of the residual
simple-root quotient give all affine primes, with $e=1$, residue degrees
equal to the factor degrees, and total residue degree $N-(r+1)$.  Thus
\[
 (r+1)+N-(r+1)=N.                                        \tag{4.12}
\]

\smallskip
\noindent\emph{Cancellation divisor $P=0$.}
The ideal $(P)$ is prime and reduced.  Since the $T$ chart is not monic
there, put $U=PT$ and
\[
 J(U,P)=\int_0^U\{P-V(Q-V)^m\}^r\,dV.
\]
The integral degree-$N$ chart
\[
 CJ(U,P)-RP^{r+1}=0                                      \tag{4.13}
\]
has the source function field: it is primitive and linear in $R$, and
$T=U/P$ over $P\ne0$.  At $P=0$ it factors as
\[
 J(U,0)=(-1)^rU^{r+1}Q^{mr}K_{m,r}(U/Q),                 \tag{4.14}
\]
where
\[
 K_{m,r}(w)=w^{-(r+1)}\int_0^w v^r(1-v)^{mr}\,dv.        \tag{4.15}
\]
The polynomial $K_{m,r}$ has degree $mr$ and is squarefree.  A nonzero
common zero of $K$ and $K'$ would make both the integral in (4.15) and
$w^r(1-w)^{mr}$ vanish; the only candidate is $w=1$, where the beta integral
is nonzero.

The generic DVR of $P=0$ has uniformizer $P$ and residue field $\C(Q,R)$.
Every branch below is unramified, with $P$ also an upstairs uniformizer:
\[
\begin{array}{c|c|c|c}
 \text{chart}&\text{residual equation}&(e,f)&\text{position}\\ \hline
 U=PS&C\int_0^S(1-tQ^m)^r\,dt-R=0&(1,r+1)&B=0\ \text{affine}\\
 U=Qw_q&K_{m,r}(w_q)=0&(1,1)&A=0\ \text{affine}\\
 U=Qw,\ w\ne w_q&K_{m,r}(w)=0&(1,1)&\text{boundary}.
\end{array}                                                \tag{4.16}
\]
The first residual polynomial is irreducible over $\C(Q,R)$ because it is
the generic equation $f(S)-R$ of degree $r+1$.  The distinguished root is
\[
 w_q=-\frac{q}{1-q},                                     \tag{4.17}
\]
where $q=h(0)$ is the selected cancellation parameter; the cancellation
identity gives $K_{m,r}(w_q)=0$.  It is precisely the affine divisor $A=0$.
There is no further affine divisorial branch: in the UFD $\C[x,y,z]$ the
equation $P=AB$ has exactly the two prime factors $A$ and $B$.  Here
$A=1+xy^m$ is primitive and linear in $x$, while
$B=A^{r+1}z+y^{m+1}h(A)$ is primitive and linear in $z$ and is not divisible
by $A$ because $h(0)=q\ne0$.  Therefore the other $mr-1$ simple roots are
boundary primes, and
\[
 (r+1)+1+(mr-1)=N.                                       \tag{4.18}
\]

\smallskip
\noindent\emph{Completed generic local rings.}
The preceding calculations can be summarized without reference to the root
charts.  If $k_Z=k(Z)$, then $\widehat{\mathcal O}_{Y,\eta_Z}=k_Z[[\pi_Z]]$,
and the completed normalization at every prime $E\mid Z$ is
$k(E)[[u_E]]$.  The four cases are:
\[
\begin{array}{c|c|c|c}
 Z&\pi_Z\mapsto\text{unit}\cdot u_E^e&k(E)&\text{position}\\ \hline
 \Delta_H&u_E^2&k_Z&\text{boundary critical prime}\\
 C=0&C&
 \begin{matrix}
  k_Z[u]/(h_2u^2-Bu+cA),\\ k_Z,\ k_Z
 \end{matrix}&
 \begin{matrix}
  W=0\text{ affine},\\ W=1\text{ affine},\ W=\rho_j\text{ boundary}
 \end{matrix}\\
 \Delta_{m,r}&u_E^{r+1}&k_Z&\text{boundary critical prime}\\
 P=0&P&
 \begin{matrix}
  k_Z[S]/(C\int_0^S(1-tQ^m)^r\,dt-R),\\ k_Z,\ k_Z
 \end{matrix}&
 \begin{matrix}
  B=0\text{ affine},\\ A=0\text{ affine},\ K_{m,r}\text{-roots boundary}.
 \end{matrix}
\end{array}                                                \tag{4.19}
\]
In the two discriminant rows, every irreducible factor of the residual
simple-root cofactor supplies an additional affine completed DVR with $e=1$
and residue field the corresponding factor field.  This accounts for every
generic DVR extension, including those which do not meet the boundary.

\smallskip
\noindent\emph{Support and absence of residual components.}
On $D(C\delta_H)$ in the weighted case and $D(P\delta_{m,r})$ in the
cancellation case, every inverse root is simple and every reconstruction
coordinate is regular.  Thus every boundary image lies in the two listed
divisors.  Equations (4.4), (4.8), (4.12), and (4.18) saturate (4.1), so no
additional normalization prime can lie over either one.  Clearing $C$ or
$P$ cannot hide a residual vertical component: its generic DVR contribution
would be an additional positive summand in a sum already equal to $N$.

The distinguished source is affine and open in the normal affine finite
normalization.  Its complement is pure of codimension one.  Indeed, at the
generic point of a hypothetical higher-codimension irreducible component,
choose an affine neighborhood avoiding the divisorial components.  Its
intersection with the affine source is affine because the ambient scheme is
separated.  Normal Hartogs extension gives the same coordinate ring on the
neighborhood and on that intersection, forcing them to be equal, a
contradiction.  A height-one prime in the finite integral cover contracts to
a nonzero height-one target prime: after excluding contraction to zero,
going down over the normal target gives equality of heights.  Hence no
higher-codimension boundary component or image is omitted.

Over a nonclosed coefficient field, factoring the residual quadratic in
(4.5) and $K_{m,r}$ in (4.16) groups geometric branches into arithmetic
primes.  The factor degrees are exactly their residue degrees, and (4.1)
again rules out a new prime after descent.  This proves the complete boundary
list and all asserted labels.  Since $mr=1$ only for $(m,r)=(1,1)$, both
second target vertices occur throughout the range $N\ge4$.

\smallskip
\noindent\emph{The full intersection diagram.}
Put $p(w)=w(1-w)^m$ and define the second degree-$mr$ polynomial
\[
 L_{m,r}(w)=w^{-(r+1)}\int_0^w\{p(w)-p(v)\}^r\,dv.       \tag{4.20}
\]
Let $S=\{P=Q=0\}$, with $R$ free.  The normalized contraction calculation
restricts the critical prime to
\[
 Q^m L_{m,r}(w)\quad\hbox{in}\quad
 \C(R)[w]/(K_{m,r})[[Q]].                               \tag{4.21}
\]
Consequently, whenever
\[
 \operatorname {Res}_w(K_{m,r},L_{m,r})\ne0,             \tag{4.22}
\]
the coefficient of $Q^m$ is a unit in the finite reduced branch algebra.
After strict henselization every geometric $K$-branch therefore intersects
the critical prime in
\[
 \C(R)[[Q]]/(Q^m).                                      \tag{4.23}
\]
One of these $mr$ branches is the affine divisor $A=0$ and the remaining
$mr-1$ are boundary primes.  Thus the critical prime has total contact
$m^2r$ with the complete $K$-cluster and contact $m(mr-1)$ with the boundary
part.  Two distinct $K$-branches force $Q=0$ and meet in the reduced stratum
$S$; adding the critical prime or any further $K$-branch does not thicken
that intersection.  These are all intersections among boundary primes.

Equation (4.22) is an exact, decidable certificate rather than a genericity
shortcut.  Under $w=-q/(1-q)$, $K_{m,r}$ is a nonzero scalar multiple of the
parameter polynomial $M_{m,r}$ after clearing $(1-q)^{mr}$.  Moreover $K$ and
$L$ have the same constant term $1/(r+1)$, whereas their linear terms are
respectively
\[
 -\frac{mr}{r+2},\qquad
 -\frac{mr(r+3)}{(r+1)(r+2)}.
\]
They are therefore not associates.  Every proved irreducibility case for
$M_{m,r}$ consequently implies (4.22).  This proves the diagram for all
$mr\le30$, the entire column $m=1$, and the other uniform arithmetic
criteria; direct resultants are additionally checked for $1\le m,r\le5$.
In particular every rung of the quadratic stationary-point ladder is covered.

There are also three uniform columns which do not use irreducibility.  Put
$M_k=K_{m,k}$, including $M_0=1$.  Scaling $v=wt$ and expanding gives
\[
 L_{m,r}=\sum_{k=0}^r(-1)^k\binom rk
 (1-w)^{m(r-k)}M_k.                                   \tag{4.24}
\]
For $r=1$, this says $L_{m,1}+K_{m,1}=(1-w)^m$; a common
root would be $w=1$, whereas
$K_{m,1}(1)=1/((m+1)(m+2))$.  In fact
\[
 \operatorname {Res}(K_{m,1},L_{m,1})
 =((m+1)(m+2))^{-m}.                                  \tag{4.25}
\]

For $r=2$, set $y=1-w$ and $z=y^m$.  The endpoint substitution
$u=1-wt$ gives
\[
 M_k=\frac1{(1-y)^{k+1}}
 \left\{\frac{k!}{\prod_{j=1}^{k+1}(mk+j)}
 -yz^k\sum_{j=0}^k\frac{(-1)^j\binom kjy^j}{mk+j+1}\right\}. \tag{4.26}
\]
At a common root, (4.24) and $M_2=0$ force $2M_1=z$, or
\[
 z\{m(m+1)y^2-2m(m+2)y+(m+1)(m+2)\}=2.               \tag{4.27}
\]
Substitution in $M_2=0$ and clearing nonzero denominators factors as
\[
 (m+1)^2(y-1)^3\{m^2y-(m+2)^2\}=0.                   \tag{4.28}
\]
The first root is $w=0$, where $K_{m,2}=1/3$.  At the second,
$y=(m+2)^2/m^2>1$, while the braced factor in (4.27) equals
$2(m+2)(5m^2+10m+4)/m^3>2$; since $z=y^m>1$, (4.27) is
impossible.  Hence (4.22) holds for every $m$ when $r=2$ as well.

For $r=3$, first note the complementary incomplete-beta identity
\[
 K_{m,r}(1-y)=\beta_r\sum_{j=0}^{mr}\binom{j+r}{r}y^j,
 \qquad
 \beta_r=\frac{r!}{\prod_{j=1}^{r+1}(mr+j)}.           \tag{4.29}
\]
The consecutive coefficient ratios are $j/(j+r)$, so the
Enestrom--Kakeya theorem puts every zero in
\[
 |y|\le \frac{mr}{mr+r}=\frac{m}{m+1}.                \tag{4.30}
\]

At a common $K_{m,3},L_{m,3}$ root, $M_3=0$, and (4.24), after division by
the nonzero $z$, becomes $z^2-3zM_1+3M_2=0$.  Put $D=1-y$ and let $T_k$
denote the sum in (4.26).  Clearing the endpoint denominators gives
\[
 az^2+bz+c=0,\qquad dz^3+e=0,                         \tag{4.31}
\]
where
\[
 \begin{aligned}
 a&=D^3+3yDT_1-3yT_2,&b&=-3\beta_1D,&c&=3\beta_2,\\
 d&=-yT_3,&&e&=\beta_3.
 \end{aligned}                                       \tag{4.32}
\]
Their resultant is
\[
 a^3e^2+3abcde-b^3de+c^3d^2.                         \tag{4.33}
\]
It factors as $(y-1)^3H_m(y)$ up to a nonzero element of $\mathbb Q(m)$, where
$H_m$ has degree six.  Put $\rho=m/(m+1)$ and
\[
 Q_m(u)=u^6H_m(\rho/u)=q_0u^6+q_1u^5+\cdots+q_6.     \tag{4.34}
\]
Let $A_m,B_m$ be the lower triangular Toeplitz matrices with entries
$(A_m)_{ij}=q_{i-j}$ and $(B_m)_{ij}=q_{6-i+j}$ for
$0\le j\le i\le5$.  After clearing a common denominator in $Q_m$, the six
leading principal minors of the Schur--Cohn matrix
\[
 A_mA_m^{\mathsf T}-B_mB_m^{\mathsf T}               \tag{4.35}
\]
are polynomials in $m$ of degrees $29,56,81,104,125,144$, respectively,
and every coefficient of every minor is strictly positive.  The matrix is
therefore positive definite for $m\ge1$.  Schur--Cohn gives
\[
 H_m(y)=0\quad\Longrightarrow\quad |y|>\frac{m}{m+1}  \tag{4.36}
\]
which is disjoint from (4.30).  Thus (4.22) also holds for every $m$ when
$r=3$.  The exact coefficient-positivity certificate is reproduced by
\path{scripts/verify_contact_resultant_endpoint_reduction.py}.
\end{proof}
